\documentclass[11pt,a4paper,french]{article}
\usepackage[T1]{fontenc}
\usepackage[utf8]{inputenc}
\usepackage{fourier}
\usepackage{xspace}
\usepackage[scaled=0.875]{helvet}
\renewcommand{\ttdefault}{lmtt}
\usepackage{amsmath,amssymb,amstext,makeidx}
\usepackage{fancybox}
\usepackage{tabularx}
\usepackage[normalem]{ulem}
\usepackage{pifont}
\usepackage[euler]{textgreek}
\usepackage{textcomp,enumitem,mathcomp}
\usepackage[table]{xcolor}
\usepackage{lscape}
\usepackage{graphicx}
\usepackage{esvect}
%tapuscrit Jean-Claude Souque
\newcommand{\euro}{\eurologo{}}
\usepackage{pst-tree,pst-plot,pst-text,pst-eucl,pst-bar,pst-func,pst-math,pstricks-add}
%\usepackage[squaren, Gray, cdot]{SIunits}
\newcommand{\R}{\mathbb{R}}
\newcommand{\N}{\mathbb{N}}
\newcommand{\D}{\mathbb{D}}
\newcommand{\Z}{\mathbb{Z}}
\newcommand{\Q}{\mathbb{Q}}
\newcommand{\C}{\mathbb{C}}
%\usepackage[left=2.1cm, right=2.1cm, top=1.9cm, bottom=2.7cm]{geometry}
\usepackage[left=3.5cm, right=3.5cm, top=2cm, bottom=3cm]{geometry}
\newcommand{\vect}[1]{\overrightarrow{\,\mathstrut#1\,}}
\renewcommand{\theenumi}{\textbf{\arabic{enumi}}}
\renewcommand{\labelenumi}{\textbf{\theenumi.}}
\renewcommand{\theenumii}{\textbf{\alph{enumii}}}
\renewcommand{\labelenumii}{\textbf{\theenumii.}}
\def\Oij{$\left(\text{O},~\vv{\imath},~\vv{\jmath}\right)$}
\def\Oijk{$\left(\text{O},~\vv{\imath},~\vv{\jmath},~\vv{k}\right)$}
\def\Ouv{$\left(\text{O},~\vv{u},~\vv{v}\right)$}
\usepackage{fancyhdr}
\makeatletter
\def\hlinewd#1{%
\noalign{\ifnum0=`}\fi\hrule \@height #1 %
\futurelet\reserved@a\@xhline}
\makeatother

\usepackage[dvips]{hyperref}
\hypersetup{%
pdfauthor = {APMEP},
pdfsubject = {Concours Général},
pdftitle = {session 2025},
allbordercolors = white,
pdfstartview=FitH}

\usepackage[np]{numprint}
\renewcommand\arraystretch{1.4}
\newcommand{\e}{\text{e}}
%\frenchbsetup{StandardLists=true}
\usepackage{babel}
\usepackage[francais]{turnthepage}
\renewcommand{\turnthepage}{\textbf{ Tournez la page S.V.P.} }
\begin{document}

\setlength\parindent{0mm}

\rhead{\textbf{A. P{}. M. E. P{}.}}
\lhead{\small { Session 2025}}
\lfoot{\small{Concours Général}}
\rfoot{}
%\renewcommand \footrulewidth{.2pt}
\pagestyle{fancy}
\thispagestyle{empty}

% \begin{center}
% {\Large\decofourleft ~ \textbf{Concours général des lycées \decofourright\\ Session 2025}}
% \end{center}

\begin{center}
CONCOURS GÉNÉRAL DES LYCÉES

\medskip

SESSION DE 2025

\medskip\

MATHÉMATIQUES

\medskip

(Classes de terminale voie générale spécialité mathématiques)

\medskip

DURÉE : 5 HEURES
\end{center}

\bigskip

\emph{L’usage de la calculatrice avec mode examen actif est autorisé.}

\emph{L’usage de la calculatrice sans mémoire \og type collège\fg{} est autorisé.}

\vspace{2cm}
\begin{center}
\fbox{
\begin{minipage}[]{0.8\textwidth}

\begin{center}
Consignes aux candidats
\end{center}

\begin{itemize}
\item  Ne pas utiliser d’encre claire
\item N’utiliser ni colle, ni agrafe
\item Ne joindre aucun brouillon
\item Ne pas composer dans la marge
\item Numéroter chaque page en bas à droite (numéro de page / nombre total de pages)
\item Sur chaque copie, renseigner l’en-tête + l’identification du concours :
\end{itemize}
\textbf{Concours / Examen : CGL \hfill Épreuve : Mathématiques \hfill Matière : MATH \hfill Session : 2025}
\end{minipage}
}
\end{center}

\vspace{2cm}

\emph{La clarté et la précision de la rédaction seront prises en compte dans l’appréciation de la copie.}

\smallskip

\emph{Le sujet comporte trois problèmes indépendants. Le candidat peut traiter les questions dans l’ordre de son choix, à condition de l’indiquer
clairement dans la copie.}

\smallskip

\emph{Dans toute question, le candidat est autorisé à utiliser les résultats de questions précédentes même s’il ne les a pas démontrés, mais à condition de le mentionner explicitement.}

\newpage
\begin{center}
\textbf{Exercice 1 : Addition sur une parabole}
\end{center}

Le plan est muni d'un repère orthonormé \Oij. Pour tout point M du plan, on note $(x_M~;~ y_M)$ ses coordonnées.

Soient $\mathcal{P}$ la parabole d'équation $y= x^2$ et $\Delta$ la droite d'équation $y =-1$.

Pour tous points $A$ et $B$ de $\mathcal{P}$ tels que $y_A\not=y_B$, on note $A\oplus B$ le point de $\mathcal{P}$ dont l'abscisse est celle du point d'intersection des droites $(AB)$ et $\Delta$.

\begin{center}
\textbf{Partie 1 : propriétés de} $\oplus$. 
\end{center}

\begin{enumerate}
\item Soient $A$ et $B$ deux points de $\mathcal{P}$ tels que $y_A\not= y_B$. Exprimer $x_{A\oplus B}$ en fonction de $x_A$ et $x_B$. 
\item Soient $A$, $B$ et $C$ trois points de $\mathcal{P}$. On suppose que
 \[y_A\not=y_B,\quad, y_{A\oplus B}\not=y_C, \quad y_B\not=y_C,\quad y_A \not= y_{B\oplus C}.\]
 Démontrer que les points $(A\oplus B)\oplus C$ et $A\oplus (B \oplus C) $ sont confondus.
 \end{enumerate}
 

Pour tout point $A$ de $\mathcal{P}$ distinct du point $O$ (de coordonnées $(0;0)$), on note $A \oplus A$ le point de $\mathcal{P}$ dont l'abscisse est celle du point d'intersection de la tangente à $\mathcal{P}$ en $A$ et de $\Delta$.
\begin{enumerate}[resume]
\item Soit $A$ un point de $\mathcal{P}$ distinct de $O$.
\begin{enumerate}
\item  Exprimer $x_{A \oplus A}$ en fonction de $x_{A}$.
\item Soit $B$ un point de $\mathcal{P}$ tel que 
\[y_{A\oplus A}\not=y_B, \quad y_A\not=y_B,\quad y_A \not=y_{A\oplus B}\]

Démontrer que les points $(A\oplus A) \oplus B$ et $A\oplus(A \oplus B)$ sont confondus.
\end{enumerate}
\end{enumerate}
\vspace{0.5cm}

\begin{center}
\textbf{Partie 2 : étude d'une suite de points}
\end{center}

Soit $A$ un point de $\mathcal{P}$. On lui assocte une suite de points $(A_n)$ définie par $A_0 = A$ et la relation de récurrence :

\[\text{pour tout entier } n\geqslant 0,\quad  A_{n+1}=\begin{cases}A_n \oplus A_n, & \text{si }A_n\not=0\\0, & \text{sinon. }\end{cases}\] 

Pour alléger les notations, on pose $x_n= x_{A_n}$. 
\begin{enumerate}[resume]
\item  On suppose, dans cette question, que $A$ est le point de coordonnées $(3;9)$.
\begin{enumerate}
\item  Démontrer que, pour tout $n$ dans $\N$, on a $x_n\not= 0$.
 \item Démontrer que la suite $(x_n)$ ne converge pas.
\end{enumerate}
\end{enumerate}

\medskip

On rappelle que, pour tout $x $ dans $Q^*$, il existe un unique couple d'entiers $(a, b)$ tels que $b \geqslant 1, \mathrm{PGCD}(a, b) = 1$ et $x= \dfrac{a}{b}$. On note alors $H(x)$ le plus grand des entiers $\vert a \vert $ et $\vert b\vert$, soit  $H(x) = \mathrm{max}\{\vert a\vert, \vert b \vert\}$.

Par exemple $H\left(\dfrac{-4}{3}\right)=4$.

 On convient de plus que $H(0)= 1$.

\smallskip

 Pour tout point $P$ de la parabole $\mathcal{P}$, tel que $x_p$ appartient à $\Q$, on pose 
\[h(P)= \ln(H(x_p)).\]

\begin{enumerate}[resume]
 \item Démontrer que, pour tout $c$ dans $\N^*$, l'ensemble $E(c)= \{P \in\mathcal{P}, x_p\in \Q \text{ et } h(P)\leqslant c \} $ est fini.
\item Soit $(a,b)\in \Z^* \times \N^*$. Démontrer que, si PGCD $(a, b) = 1$, alors 

PGCD $ (a^2-b^2,  ab) =1$.
\item Démontrer qu'il existe deux réels $m$ et $M$ tels que pour tout point $P$ de la parabole $\mathcal{P}$ tel que $x_p$
appartient à $\Q^*$, on a
\[m+h(P\oplus P)\leqslant 2h(P)\leqslant h(P\oplus P)+M.\]

\item Soit $(u_n)$ une suite réelle. On pose

\[v_0=0 \ \text{ et, pour tout } n\geqslant 1,\quad v_n=\sum_{k=0}^{n-1} \vert u_{k+1}-u_k\vert.\]


On suppose que la suite $(v_n)$ est majorée et on souhaite démontrer que la suite $(u_n)$ converge :
\begin{enumerate}
\item  Démontrer que la suite $(v_n)$ converge.
\item Démontrer que pour tout entier naturel $n$, on a
\[0\leqslant u_{n+1}- u_n + \vert u_{n+1} - u_n\vert\leqslant  2\vert u_{n+1} - u_n\vert .\]
\item En déduire que la suite $(u_n+ v_n)$ converge. 
\item Conclure.
\end{enumerate}
\item Soit à nouveau le point $A$ de coordonnées (3;9) et soit $(A_n)$ la suite récurrente associée à $A$. Pour
tout entier naturel $n$, on pose $ t_n=\frac{h(A_n)}{2^n}$. Démontrer que la suite $(t_n)$ converge.
\end{enumerate}

\bigskip

\begin{center}
\textbf{Exercice 2 : Suite positive et suite bornée}
\end{center}


Pour tout réel $\alpha$, on appelle suite associée à $\alpha$ la suite $(u_n)$ définie par $u_0=u_1 = 1 $ et
 \[u_{n+2}=u^2_{n+1}-\alpha u_n^4\]
pour tout entier $n \geqslant 0$. On dit que  $\alpha$ \emph{vérifie la propriété} $\mathcal{P}$ si tous les termes de la suite $(u_n)$ assoctée à $\alpha$ sont strictement positifs, et que $\alpha$ \emph{vérifie la propriété} $\mathcal{B}$ si la suite $(u_n)$ associée à $\alpha$ est bornée.

\medskip
\begin{center}
\textbf{Partie 1 : Propriété $\mathcal{P}$}
\end{center}

\medskip

\begin{enumerate}
\item  Quels sont les réels $ \alpha $ qui vérifient la propriété $\mathcal{P}$ et qui appartiennent 
\begin{enumerate}
\item  à l'intervalle $[1~;~+\infty[$ ? 
\item à l'intervalle $]-\infty~;~0]$?
\end{enumerate}
\item Soit $\alpha$ un réel appartenant à l'intervalle $]0; 1[$ et $(u_n)$ la suite qui lui est associée; on suppose, dans cette question, que $\alpha$ vérifie la propriété $\mathcal{P}$.
\begin{enumerate}
\item  Démontrer que $0<u_{n+1}\leqslant u_n\leqslant 1$ pour tout entier $n \geqslant 0$.
\item Quelle est la limite de la suite $(u_n)$?
\item Pour tout réel $n \geqslant 0$, on pose $x_n=\dfrac{u_{n+1}}{u_n^2}$. Exprimer $x_{n+1}$ en fonction de $\alpha$ et de $x_n$.
\item Démontrer que la suite $(x_n)$ admet une limite finie, que l'on notera $x_{\infty}$ et exprimer $x_{\infty}^2 (1- x_{\infty})$ en fonction de $\alpha$.
\item En déduire que $\alpha\leqslant \dfrac{4}{27}$.
\end{enumerate}
\item  Quels sont les réels $\alpha$ qui vérifient la propriété $\mathcal{P}$? 
\end{enumerate}

\vspace{0.5cm}

\begin{center}
Partie 2 : Propriété $\mathcal{B}$
\end{center}

\medskip

\begin{enumerate}[resume]
\item  Quels sont les réels $\alpha$ qui vérifient la propriété $\mathcal{B}$ et qui appartiennent
\begin{enumerate}
\item  à l'intervalle $]-\infty~;~0[$? 
\item à l'intervalle [0:1]? 
\end{enumerate}
\item Soit $\alpha$ un réel appartenant à l'intervalle $]2~;~+\infty[$ et $(u_n)$ la suite qui lui est associée. Pour tout $n\geqslant 0$, on note $v_n$, le plus grand des réels $\vert u_0\vert$, $\vert u_1\vert,\dots,\vert u_n\vert$, soit 

$v_n = \mathrm{max}\{\vert u_0\vert,\vert u_1\vert,\dots,\vert u_n\vert\}$. 
\begin{enumerate}
\item  Démontrer, pour tout entier $n \geqslant 1$, que $v_n=\vert u_{n-1}\vert$, ou $v_n=\vert u_n\vert$. 
\item Le réel $\alpha$ vérifie-t-il la propriété $\mathcal{B}$?
\end{enumerate}
\item Soit $\alpha$ un réel appartenant à l'intervalle $\left[\dfrac{9}{7}~;~2\right]$ et $(u_n)$ la suite qui lui est associée.
\begin{enumerate}
\item  Pour tout réel $x$, on pose $ P(x)= ax^3- x -1$ et $Q(x) =8x^3 -116x^2 + 494x- 441$. Démontrer que

\[P\left(\dfrac{11-2\alpha}{7}\right)=\dfrac{(2-\alpha)Q(\alpha)}{7^3}\]

\item Étudier les variations de la fonction $x\mapsto P(x)$ sur l'intervalle $[1~;~+\infty[$ et de la fonction $x\mapsto Q(x)$ sur l'intervalle $[1~;~2]$.
\item Quel est le signe de $Q(\alpha)$?
\item Comparer les réels $\dfrac{11-2\alpha}{7}$ et $-u_3$.
\item Le réel $\alpha$ vérifie-t-il la propriété $\mathcal{B}$ ?
\end{enumerate}
\item Soit $\alpha$ un réel appartenant à l'intervalle $\left]1~;~\frac{9}{7}\right[$. Pour tout réel $x$, on pose $S_0(x)= x$, $S_1(x) = 1$, puis
\[S_{k+2}(x)=S_{k+1}(x)^2-\alpha S_k(x)^4\]

 pour tout entier $k\geqslant 0$. On \emph{admet} pour l'instant qu'il existe un réel $\mathrm{t}(\alpha) > 1$ tel que 

\[1< S_4(x)\leqslant S_3(x)^2\leqslant \mathrm{t}(\alpha)^2S_4(x)\]

pour tout réel $x \in [1~;~\mathrm{t}(\alpha)]$.
\begin{enumerate}
\item  Soit $n \geqslant 0$ un entier tel que $1\leqslant u_{n+1}\leqslant u_n^2<\mathrm{t}(\alpha)^2u_{n+1}$ ; On pose $x_n=\sqrt{\dfrac{u_n^2}{u_{n+1}}}$. 

Démontrer que $ u_{n+k}=u_{n+1}^{2^{k-1}}S_k(x_n)  $
pour tout entier $k\geqslant 1$.
 \item Le réel  $\alpha$ vérifie-t-il la propriété $\mathcal{B}$ ?
 \end{enumerate}
\item On démontre maintenant que l'unique réel positif $\mathrm{t}(\alpha)$ tel que $\mathrm{t}(\alpha)^4 = \dfrac{\alpha+2}{3}$ vérifie les hypothèses
de la question \textbf{7.} 
\begin{enumerate}
\item  Démontrer que $S_4(x)\leqslant S_3(x)^2$ pour tout réel $x$.
\item Démontrer que $0< S_2(x)^2\leqslant\dfrac{7(\alpha-1)}{12}$ pour tout réel $x\in[1~;~\mathrm{t}(\alpha)] $
\item Démontrer que $1< S_4(x)$ pour tout réel $x\in[1;\mathrm{t}(\alpha)]$.
\item Démontrer que $S_3(x)^2\leqslant \dfrac{2+x^4}{3}S_4(x)\leqslant x^2S_4(x)$ lorsque $x =\mathrm{t}(\alpha)$.
\item Étudier les variations de la fonction $x \mapsto \dfrac{S_4(x)}{S_3(x)^2}$ sur l'intervalle $[1~;~\mathrm{t}(\alpha)]$.
\item Conclure.
\end{enumerate}
 \item Quels sont les réels $\alpha$ qui vérifient la propriété $\mathcal{B}$ ?
 
\end{enumerate}

\vspace{1cm}

\begin{center}
\textbf{Exercice 3}
\end{center}

On dit qu'une fonction $f$ vérifie la propriété $\mathcal{E}$ si, pour tout réel $x$, 
\[f(x+1)= \dfrac{1}{2}+\sqrt{f(x)-f(x)^2}.\]
\begin{enumerate}
\item  Proposer une fonction continue qui vérifie la propriété $\mathcal{E}$.
\item On rappelle qu'une fonction $g$ définie sur $\R$ est périodique s'il existe un réel $T > 0$ tel que
\[\text{pour tout réel } x,~ g(x+T) = g(x).\]
Soit $f$ une fonction vérifiant $\mathcal{E}$. Démontrer que $f$ est périodique. 
\item Proposer une infinité de fonctions continues $f$ vérifiant $\mathcal{E}$ et $f(0) =\dfrac{1}{2}$.

\end{enumerate}




\end{document}
